% paper/sections/09d_pressure_summary.tex
%
% §9 章末まとめ：精度バランスと設計上のトレードオフ
% ※ 旧所在: 09_ccd_poisson.tex 内 → §9 章末に移動
%   理由：HFE（§9.4）・BC（§9.6）を全て読んだ後に精度全体像を示す構成に変更
%% ═══════════════════════════════════════════════════════════════════════════════
\subsection{本章のまとめ：精度バランスと設計上のトレードオフ}
\label{sec:pressure_accuracy_summary}
%% ═══════════════════════════════════════════════════════════════════════════════

変密度 PPE の導出（§~\ref{sec:projection_derivation}）は第~\ref{sec:collocate}章で行った．
本章では，CCD による $\Ord{h^6}$ 離散化（§~\ref{sec:ccd_poisson_matrix}），
分相 PPE 解法（§~\ref{sec:split_ppe_framework}），
HFE による界面越し場延長（§~\ref{sec:field_extension}），
各相内の欠陥補正法による求解（§~\ref{sec:defect_correction_main}），
境界条件と数値安定性（§~\ref{sec:ppe_bc_stability}）を論じた．
以下に，各コンポーネントの精度を整理する．

\begin{table}[htbp]
\centering
\caption{数値手法コンポーネント別の精度まとめ}
\label{tab:accuracy_summary}
\renewcommand{\arraystretch}{1.3}
\footnotesize
\begin{tabular}{lp{0.28\linewidth}cc}
\toprule
コンポーネント & 手法 & 空間精度 & 時間精度 \\
\midrule
CLS 移流（界面追跡） & Dissipative CCD + TVD-RK3 & $\Ord{h^5}$§ & $\Ord{\Delta t^3}$ \\
幾何量計算（曲率等） & CCD 法 & $\Ord{h^6}$ & — \\
PPE（離散化＋求解） & 分相 PPE + DC $k{=}3$ + CCD 勾配（Balanced-Force） & $\Ord{h^5}$/$\Ord{h^6}$‡ & — \\
NS 粘性項 & CCD + Crank--Nicolson（陽的交差項 $\Ord{h^6}$；ADI 陰的対角項 $\Ord{h^2}$）${}^{\P}$ & $\Ord{h^2}$/$\Ord{h^6}$${}^{\P}$ & $\Ord{\Delta t^2}$ \\
NS 対流項（Predictor） & CCD（1次微分）+ AB2 外挿 & $\Ord{h^6}$ & $\Ord{\Delta t^2}$ \\
Projection スプリッティング & IPC 法（van Kan 1986） & — & $\Ord{\Delta t^2}$ \\
HFE（分相 PPE 基盤，§~\ref{sec:field_extension}） & 最近接点 Hermite 補間 & $\Ord{h^6}$ & — \\
\midrule
\textbf{全体（空間律速：CSF モデル誤差）} & & $\boldsymbol{\Ord{h^2}}$† & — \\
\textbf{全体（時間律速：AB2+IPC 整合）} & & — & $\boldsymbol{\Ord{\Delta t^2}}$ \\
\bottomrule
\end{tabular}

\smallskip
\noindent†\ CSF 正則化 $\delta_\varepsilon$（$\varepsilon = \Ord{h}$）によるモデル誤差 $\Ord{\varepsilon^2} = \Ord{h^2}$ が
全体の空間精度を律速する．\\
\noindent‡\ 分相 PPE + DC $k{\geq}3$ では各相内 $\Ord{h^7}$（Dirichlet 超収束）/ $\Ord{h^5}$（Neumann BC 律速）を全密度比で達成（§~\ref{sec:verify_split_ppe_dc}）；速度補正の勾配は CCD $\Ord{h^6}$．\\
\noindent§\ 単位ステップ質量誤差の大域積分推定（§~\ref{sec:dccd_conservation}）；
局所切断誤差は $\Ord{h^6}$（CCD）から $\Ord{\varepsilon_d h^2}$（DCCD フィルタ）が加わった実効精度．\\
\noindent¶\ ADI 陰的スウィープの対角係数には2次 FD（$\Ord{h^2}$）を使用；交差微分項は CCD $\Ord{h^6}$ で陽的に評価（§~\ref{sec:algorithm}）．
$Re \gg 1$ では粘性項の寄与 $\sim \Ord{h^2/Re}$ が CSF モデル誤差 $\Ord{h^2}$ と同次以下に収まるため実効精度への影響は小さい．
\end{table}

\begin{tcolorbox}[warnbox, title={精度ミスマッチと設計上のトレードオフ}]
\begin{enumerate}
  \item \textbf{空間精度：}
        全数値コンポーネントが $\Ord{h^5}$--$\Ord{h^6}$ に揃っている（平滑領域）．
        全体精度の律速は CSF 正則化 $\delta_\varepsilon$ による $\Ord{h^2}$ モデル誤差である．
        分相 PPE + DC $k{=}3$ は各相内 $\Ord{h^5}$（Neumann BC 律速）を全密度比で達成し，
        HFE（§~\ref{sec:field_extension}）が界面ジャンプ条件を $\Ord{h^6}$ で課す．

  \item \textbf{時間精度：}
        AB2 + IPC により全体 $\Ord{\Delta t^2}$（§~\ref{sec:ipc_derivation}）．
        対流（AB2），粘性（Crank--Nicolson），Projection（IPC）の3要素が統一されている．
\end{enumerate}
\end{tcolorbox}

\begin{tcolorbox}[mybox, title={設計選択の根拠：分相 PPE + DC + CCD 勾配の統合設計}]
\label{box:ppe_design_rationale}
本稿の PPE 求解戦略は\textbf{分相 PPE 解法}を主軸とし，以下の3つの原則から帰結する．

\begin{enumerate}
  \item \textbf{分相 PPE による密度比非依存化}．
        各相で定数密度の PPE を独立に解き，界面ジャンプ条件 $[p] = \sigma\kappa$ を
        HFE（§~\ref{sec:field_extension}）で課す（§~\ref{sec:split_ppe_framework}）．
        界面を横切るステンシルが完全に排除されるため，
        DC $k{=}3$ は各相内で $\Ord{h^7}$（Dirichlet）/ $\Ord{h^5}$（Neumann BC 律速）を
        \textbf{全密度比}（$\rho_l/\rho_g = 1$--$1000$）で達成する
        （§~\ref{sec:verify_split_ppe_dc}）．

  \item \textbf{Balanced-Force 整合性：同一演算子の強制}．
        Balanced-Force 条件は，圧力勾配 $\bnabla p^{n+1}$ と表面張力 $\sigma\kappa\bnabla\psi$ に
        \textbf{同一の離散微分演算子}を使用することで，静止液滴での寄生流れを離散レベルで消去する
        （§~\ref{sec:rc_balanced_force}）．
        速度補正 $\delta\bu = -\Delta t\,\bnabla p^{n+1}/\rho$ の勾配に CCD $\Ord{h^6}$ を使い，
        表面張力項 $\kappa\bnabla\psi$ にも同一の CCD 演算子を適用することで，
        この整合性を $\Ord{h^6}$ 精度で保証する．

  \item \textbf{直接法による残差ゼロ保証}．
        DC 法の補正方程式 $L_L$（§~\ref{sec:dc_sweep}）を LU 分解（直接法）で厳密に解くことで，
        DC 反復の各ステップで $L_L$ の離散残差は機械精度まで消去される．
        直接法はこの系統誤差を原理的に排除する（§~\ref{sec:defect_correction_main}参照）．
\end{enumerate}

\noindent\textbf{全体精度の律速：}
CSF 正則化 $\delta_\varepsilon$（$\varepsilon=\Ord{h}$）によるモデル誤差
$\Ord{\varepsilon^2}=\Ord{h^2}$ が全体の空間精度を律速する（表~\ref{tab:accuracy_summary}†）．
分相 PPE + DC $k{=}3$ の $\Ord{h^5}$ は CSF 誤差を十分に上回り，
PPE ソルバが精度ボトルネックとならないことを保証する．
\end{tcolorbox}
